%!TEX root = ../main.tex
\section{Scope and Review Protocol}
\label{sec:method}

We conduct a systematic scoping review rather than a statistical meta-analysis.
The literature does not share one data intervention, learner, dataset, or outcome protocol, so success-rate differences cannot be pooled into a meaningful estimate.
The purpose of this section is to establish which works support the survey and how strongly their experiments support claims about data value.
Detailed search strings, update records, and coding fields are reported in Appendix~\ref{app:protocol}.

\subsection{Scope and Eligibility}

The four taxonomy questions introduced in Section~\ref{sec:introduction} also serve as the review questions.
They cover the operation applied to data, the signal used to choose it, the evidence used to validate it, and the point in the lifecycle at which it changes.
This section therefore defines eligibility and evidence rules rather than introducing a second question set.

A work is eligible when it describes an identifiable operation on embodied data, a source design that changes downstream compatibility or coverage, or an empirical finding about data composition.
Peer-reviewed papers, arXiv papers, dataset and benchmark papers, and sufficiently detailed technical reports are eligible.
Official project documentation is used only for release, format, and implementation facts.
We exclude work that only consumes an existing dataset, changes only the policy architecture, or presents a simulator without a transferable data operation.
Navigation and driving are outside the core scope unless a preparation mechanism transfers directly to embodied policy learning.

\subsection{Search and Screening}

The search covers January 2018 through August 2, 2026.
We queried DBLP and arXiv for combinations of robotics terms with curation, alignment, labeling, selection, generation, and feedback-driven training.
Backward and forward citation tracing and checks of major datasets, laboratories, and industrial reports supplemented the database search.
A targeted update on August 2 added recent industrial data recipes, robot post-training systems, data flywheels, and relevant language-model data-engineering work.

The resolved database search produced 1,160 title-deduplicated candidates.
Title and abstract screening removed records outside the survey boundary, and primary-text screening required a recoverable operation, decision signal, or finding about data composition.
Duplicate preprint and published versions were merged.
The final evidence map contains 211 records.
Technical analysis draws only on the 177 works for which we examined the primary-text sections needed to verify each coded method or evidence claim.
The remaining 34 records form an update index and do not support technical claims.
Figure~\ref{fig:review-flow} summarizes this distinction between discovery coverage and the evidence used in synthesis.

\begin{figure*}[t]
\centering
\begin{tikzpicture}[
  node distance=3.5mm,
  box/.style={draw, rounded corners=2pt, align=center, text width=14.5cm,
    inner sep=3.5pt, fill=gray!6},
  branch/.style={draw, rounded corners=2pt, align=center, text width=6.8cm,
    inner sep=3.5pt, fill=blue!6},
  arr/.style={-{Latex[length=1.8mm]}, thick, draw=black!65}
]
\node[box] (sources) {28 resolved DBLP and arXiv queries: 1,160 title-deduplicated discovery records\\$+$ backward and forward citation tracing $+$ major dataset, benchmark, lab, and system reports};
\node[box, below=of sources] (map) {\textbf{Evidence map: 211 unique works}\\2018--2026 baseline search $+$ targeted industrial and closed-loop update};
\node[branch, below=5mm of map, xshift=-3.75cm] (deep) {177 core works\\examined in primary text};
\node[branch, below=5mm of map, xshift=3.75cm] (abstract) {34 update-index records\\not used for technical claims};
\node[box, below=15mm of map] (coding) {Core synthesis: operations, decision signals, training roles, validation evidence, lifecycle stage, assumptions, and limitations};
\draw[arr] (sources) -- (map);
\draw[arr] (map) -- (deep);
\draw[arr] (map) -- (abstract);
\draw[arr] (deep) -- (coding);
\end{tikzpicture}
\caption{Auditable evidence-flow summary. Update-index records establish coverage but do not support technical claims until their primary text is verified. Appendix~\ref{app:protocol} provides the queries and criteria.}
\label{fig:review-flow}
\end{figure*}


\subsection{Coding and Evidence}

For each of the 177 core works, we record the data operation, the object it changes, the decision signal, and the reported validation.
We also record the intended training role, human involvement, cost, generalization setting, failure modes, provenance, and release status when the source reports them.
Composite systems receive one code for each material operation.
This preserves the difference between labeling, filtering, generation, and assigning a training role even when a paper calls the entire pipeline ``curation.''

Evidence is separated into three levels because these levels support different conclusions.
Data-level evidence shows that an intended property changed, such as a reduction in missing records or duplicates.
Proxy evidence shows agreement with a model, simulator, or estimator.
Downstream evidence trains a policy under a controlled budget and tests the resulting behavior in simulation or on a robot.
Only the third level directly establishes a policy benefit, although it may still be narrow in task, learner, or embodiment.

Peer-reviewed experiments and reproducible primary papers provide the strongest support for causal claims.
Preprints can provide method and experimental evidence with greater uncertainty about review status.
Technical reports are often the only source for industrial recipes, but end-to-end system results do not isolate the contribution of each preparation step.
We mark cross-paper interpretations as synthesis rather than attributing them to a single source.

\subsection{Limitations}

Public evidence under-represents private industrial operations and negative results.
Terminology is unstable, especially around data engines, updates after initial training, models of environment evolution, and self-improvement.
Recent 2026 papers are also dominated by preprints.
Multiple search routes and the public evidence map reduce the risk of omission but do not remove it.
We therefore avoid cross-paper rankings, preserve version-specific claims, and state when a conclusion depends on evidence from only one model, task, or robot.
